\section{Conclusion}\label{sec:conclusion}

We have established a rigorous mathematical framework for analyzing
chain-of-thought reasoning as a discrete dynamical system on the probability
simplex. Our main contributions are:

\begin{enumerate}[label=(\arabic*)]
\item Exponential convergence in $W_1$ under Dobrushin contraction
  (Theorem~\ref{thm:exponential}), with rate governed by the Wasserstein
  contraction coefficient $\eta(P) \le \delta(P)$.

\item A spectral characterization of convergence rates
  (Theorem~\ref{thm:spectral}): the $W_2$ rate is $(1 - \gamma)^{1/2}$,
  verified empirically with $r = 0.979$ correlation across 30 random
  transition matrices.

\item An exact spectral gap formula $\gamma = K\varepsilon/(K-1)$ for
  metastable chains (Theorem~\ref{thm:metastable}), matching numerical
  computation to machine precision.

\item A phase transition in convergence behavior at critical coupling
  $\varepsilon_c = (K-1)\gamma_{\mathrm{intra}}/K$
  (Theorem~\ref{thm:phase}), separating slow-mixing and fast-mixing
  regimes.
\end{enumerate}

The framework provides a principled basis for understanding when and why
iterative reasoning converges, with explicit dependence on problem structure
(number of clusters $K$, coupling strength $\varepsilon$) and spectral
properties of the transition operator. We expect the tools developed
here---particularly the connection between metastable structure and spectral
gap---to be applicable more broadly to the analysis of iterative inference
procedures in machine learning.
